% =====================================================================
% Figure contract
% 核心结论：CUA 产业由基础设施、模型与 API、产品三层向上承托，治理横跨上层，
%           而部署与行业成熟度共同决定研究结论能否落地。
% 图原型：mode A / 分层架构；沿用 cua-survey-fig1-overview.tex 的低饱和配色，
%           并采用 layer-stack.tex 的整宽层条 + 内嵌卡片构造。
% 证据层级：hero = L1--L3 三层上升结构；支撑 = 右侧 L4 治理栏；收束 = 底部注释条。
% 能不能更少？：不能；三层、两条上行关系、治理侧栏及部署注释均为指定语义。
%
% Step ① 设计文档（Form A / 中等复杂度分层架构）
% 0 参考：复用既有 CUA 图的 muted palette 与字体节奏，不复刻外部位图。
% 1 领域：Computer-Use Agents 中文综述；agent、API、RPA、VM 保留英文。
% 2 格式：TikZ standalone；XeLaTeX + fontspec + xeCJK。
% 3 画布：W_ink 约 18.4 cm，W_print = 18.4 cm；目标印刷最小 CJK 字号 7.5 pt。
% 4 模块：主体三条 13.8×2.3 cm 横条；L1/L2 各两卡，L3 三卡；右侧 L4 两卡竖排。
% 5 连线：L1→L3 左侧正交 rail；L2→L3 右侧直线；L4→L3 短虚线。
% 6 空间：主层 x=0--13.8；治理栏 x=14.55--17.95；左 rail x=-0.45。
% 7 强调：四层四色；层条 1.0 pt 描边，内卡 0.65 pt，治理栏虚线 1.1 pt。
% 8 密度：中等档；纯结构语义无数值，不嵌入 chart/heatmap。
% 9 ASCII 草图：
%       ┌────────────── L3 产品（三卡） ──────────────┐  ┌─ L4 治理 ─┐
%   ↑   └────────────────────────────────────────────┘←┆ 凭证与权限 ┆
%   │          ↑                                      ┆ 可观测审计 ┆
%   │   ┌──────── L2 模型与 API（两卡） ─────────────┐  └───────────┘
%   │   └────────────────────────────────────────────┘
%   └───┌────────── L1 基础设施（两卡） ─────────────┐
%       └────────────────────────────────────────────┘
%       ┌────────────── 部署缺口注释条 ─────────────────────────────┐
%       └───────────────────────────────────────────────────────────┘
% 10 可视化决策：输入无量化数据，结构本身即视觉主角；不添加无来源信息 panel。
%
% Pre-flight：P1 设计文档已写；P2 三层共享 x/宽度；P3 三条边均列于代码；
% P4 N/A；P5 左/右 corridor 分离；P6 内卡均留足 parent padding；P7 已读 lessons。
% =====================================================================

\documentclass[border=8pt]{standalone}
\usepackage{fontspec}
\usepackage{xeCJK}
\usepackage{xcolor}
\usepackage{tikz}
\setCJKmainfont{Heiti SC}
\setCJKsansfont{Heiti SC}
\usetikzlibrary{arrows.meta,bending,positioning,calc,backgrounds,fit}

% Low-saturation academic palette, aligned with the survey overview figure.
\definecolor{ink}{HTML}{27364A}
\definecolor{muted}{HTML}{748092}
\definecolor{blueLine}{HTML}{6D879E}
\definecolor{blueFill}{HTML}{E7EEF4}
\definecolor{tealLine}{HTML}{648D8A}
\definecolor{tealFill}{HTML}{E5F0EE}
\definecolor{purpleLine}{HTML}{817594}
\definecolor{purpleFill}{HTML}{EEEAF2}
\definecolor{roseLine}{HTML}{9A777B}
\definecolor{roseFill}{HTML}{F3E9EA}
\definecolor{greyLine}{HTML}{8A929D}
\definecolor{greyFill}{HTML}{F1F3F5}
\definecolor{flowLine}{HTML}{53697C}

\pgfdeclarelayer{background}
\pgfsetlayers{background,main}

\tikzset{
  every node/.style={
    font=\sffamily\fontsize{8.8}{10.6}\selectfont,
    text=ink
  },
  layerbar/.style={
    rectangle,
    rounded corners=6pt,
    line width=1.0pt,
    minimum width=13.80cm,
    minimum height=2.30cm,
    inner sep=0pt
  },
  layerlabel/.style={
    align=left,
    text width=3.30cm,
    font=\sffamily\fontsize{9.3}{11.2}\selectfont\bfseries
  },
  innercard/.style={
    rectangle,
    rounded corners=4pt,
    line width=0.65pt,
    minimum height=1.22cm,
    align=center,
    inner sep=5pt,
    fill=white
  },
  twocard/.style={
    innercard,
    minimum width=4.10cm,
    text width=3.65cm
  },
  threecard/.style={
    innercard,
    minimum width=2.65cm,
    text width=2.28cm
  },
  governancecard/.style={
    innercard,
    minimum width=2.72cm,
    minimum height=1.04cm,
    text width=2.30cm
  },
  governancezone/.style={
    rectangle,
    rounded corners=7pt,
    draw=roseLine,
    fill=roseFill,
    densely dashed,
    line width=1.1pt,
    inner sep=0pt
  },
  summarybar/.style={
    rectangle,
    rounded corners=4pt,
    draw=greyLine!70,
    fill=greyFill,
    line width=0.65pt,
    minimum width=17.95cm,
    minimum height=1.15cm,
    text width=16.95cm,
    align=center,
    inner sep=5pt,
    font=\sffamily\fontsize{8.5}{10.4}\selectfont
  },
  arrow/.style={
    -{Stealth[length=5pt 1.5,width'=0pt 0.6,sep=0pt 0.5]},
    line width=1.0pt,
    shorten >=2pt,
    shorten <=1pt,
    color=flowLine
  },
  arrow short/.style={
    -{Stealth[length=3pt 1.0,width'=0pt 0.5,sep=0pt 0.3]},
    line width=0.8pt,
    shorten >=1pt,
    shorten <=0pt,
    color=flowLine
  },
  governance arrow/.style={
    -{Stealth[length=3pt 1.0,width'=0pt 0.5,sep=0pt 0.3]},
    line width=0.8pt,
    densely dashed,
    shorten >=1pt,
    shorten <=0pt,
    color=roseLine
  }
}

\begin{document}
\sffamily
\begin{tikzpicture}

% Title.
\draw[line width=2.2pt,color=flowLine] (0,9.42) -- (0,10.02);
\node[
  anchor=west,
  font=\sffamily\fontsize{11.2}{13.4}\selectfont\bfseries,
  text=flowLine
] at (0.25,9.72) {产业格局:四层结构(§9)};

% Main three-layer stack: identical width and strict vertical alignment.
\node[layerbar,fill=blueFill,draw=blueLine] (l1) at (6.90,1.40) {};
\node[layerbar,fill=tealFill,draw=tealLine] (l2) at (6.90,4.55) {};
\node[layerbar,fill=purpleFill,draw=purpleLine] (l3) at (6.90,7.70) {};

% Layer labels.
\node[layerlabel,anchor=west,text=blueLine] at ([xshift=0.45cm]l1.west)
  {L1 基础设施 §9.6};
\node[layerlabel,anchor=west,text=tealLine] at ([xshift=0.45cm]l2.west)
  {L2 模型与 API\\§9.1·9.4};
\node[layerlabel,anchor=west,text=purpleLine] at ([xshift=0.45cm]l3.west)
  {L3 产品 §9.2–9.5};

% Child cards stay comfortably inside each layer boundary.
\node[twocard,draw=blueLine] (cloudBrowser) at (6.10,1.40) {云端浏览器};
\node[twocard,draw=blueLine] (sandboxVm) at (10.80,1.40) {沙箱与桌面 VM};

\node[twocard,draw=tealLine] (closedApi) at (6.10,4.55) {闭源 CU API};
\node[twocard,draw=tealLine] (openWeights) at (10.80,4.55) {开源与端侧权重};

\node[threecard,draw=purpleLine] (consumerAgent) at (5.55,7.70) {消费级 agent};
\node[threecard,draw=purpleLine] (enterpriseRpa) at (8.55,7.70)
  {企业自动化\\与 RPA};
\node[threecard,draw=purpleLine] (verticalAgent) at (11.55,7.70) {垂直 agent};

% Governance sidebar: fit-by-construction over the exact L2+L3 vertical span.
\coordinate (govNW) at (14.55,8.85);
\coordinate (govSE) at (17.95,3.40);
\node[
  anchor=north,
  font=\sffamily\fontsize{9.2}{11.0}\selectfont\bfseries,
  text=roseLine
] (govTitle) at (16.25,8.46) {L4 治理 §9.7–9.8};
\node[governancecard,draw=roseLine] (credentials) at (16.25,6.92) {凭证与权限};
\node[governancecard,draw=roseLine] (observability) at (16.25,5.25) {可观测与审计};
\begin{pgfonlayer}{background}
  \node[governancezone,fit=(govNW)(govSE)(govTitle)(credentials)(observability)] (governance) {};
\end{pgfonlayer}

% Solid upward relationships, intentionally split across left/right corridors.
\coordinate (l1Source) at (l1.west);
\coordinate (leftRailLow) at (-0.45,1.40);
\coordinate (leftRailHigh) at (-0.45,6.03);
\coordinate (l3LeftShelf) at (1.25,6.03);
\coordinate (l3LeftTarget) at ([xshift=1.25cm]l3.south west);
\draw[arrow short,rounded corners=5pt]
  (l1Source) -- (leftRailLow) -- (leftRailHigh)
  -- (l3LeftShelf) -- (l3LeftTarget);

\coordinate (l2RightSource) at ([xshift=-1.15cm]l2.north east);
\coordinate (l3RightTarget) at ([xshift=-1.15cm]l3.south east);
\draw[arrow short] (l2RightSource) -- (l3RightTarget);

% Governance influences the product layer through a distinct dashed link.
\draw[governance arrow] (governance.west |- l3.east) -- (l3.east);

% Full-width deployment note.
\node[summarybar] (deploymentNote) at (8.975,-0.75)
  {部署缺口:成本·延迟·业务价值(§9.9)与行业成熟度(§9.10)决定研究结论能否落地};

\end{tikzpicture}
\end{document}
